\documentclass[11pt]{article}

\usepackage[margin=1in]{geometry}
\usepackage{times}
\usepackage{graphicx}
\usepackage{booktabs}
\usepackage{amsmath}
\usepackage{hyperref}
\usepackage{caption}
\usepackage{enumitem}
\usepackage{natbib}

\title{Drowning in Meaning: Steering Figurative-vs-Literal\\ Interpretation Beyond Idioms}
\author{Itamar Vasserman (ID 205949985)}
\date{20/08/2026}

\begin{document}
\maketitle
\vspace{-2em}
\begin{center}
Submitted as final project report for the NLP course, RU, 2026
\end{center}
\vspace{1em}

\section{Introduction}

Idioms such as \emph{``take a bow''} or \emph{``green light''} admit two readings: a
figurative, conventionalized meaning, and a literal, compositional one. \citet{idiosteer}
(\emph{Steering Idioms}, IdioSteer) recently showed that this ambiguity can be controlled
mechanistically: a single mean-difference vector, extracted from an LLM's residual stream
and added back in at generation time, reliably pushes a model's continuation of an
ambiguous idiom prefix from a figurative reading toward a literal one (and vice versa),
while largely preserving fluency.

That result raises a natural question, and the one this project addresses: \textbf{is the
direction IdioSteer discovers a general figurative-vs-literal axis in the model's
representation space, or is it specific to the closed, conventionalized class of
expressions we call idioms?} Idioms are a special case of figurative language --
frozen, non-compositional multiword expressions with a fixed conventional meaning.
Metaphors and similes are figurative in the same broad sense (a literal reading exists
alongside an intended non-literal one) but are typically novel, compositional, and
constructed on the fly rather than memorized as a unit. If the steering vector encodes a
general semantic distinction between ``interpret this compositionally'' and ``interpret
this as a fixed non-literal unit,'' we would expect it to transfer to metaphors and
similes as well. If it instead encodes something closer to ``this specific memorized
idiom has a non-literal sense,'' we would expect little or no transfer.

We test this by (1) reproducing the IdioSteer pipeline -- steering vector construction,
cross-layer additive injection, and LLM-judge evaluation -- on a base LLM and an idiom
benchmark, and (2) applying the resulting, idiom-derived steering vector, \emph{unmodified},
to a new hand-built benchmark of 100 non-idiomatic figurative prefixes (conventional
metaphors, novel metaphors, and similes) with the same ambiguous-prefix structure. We
compare the figurative/literal rate shift induced by the same vector, at the same
coefficients, across idioms and each figurative category. We do not train a second,
figurative-derived steering vector for geometric comparison in this pass (IdioSteer's
optional Step 4); that is left as future work (Section~\ref{sec:conclusion}).

Our main finding is a \textbf{partial generalization}: the idiom vector shifts literal
rate in the expected direction for similes and conventional metaphors, with a smaller
effect size than on idioms, and shows no effect in that direction for novel metaphors --
though the latter category's baseline is already literal-skewed, leaving little room for
a further literal push. We also find that our own idiom replication is directionally
consistent with, but numerically much weaker than, IdioSteer's reported result, which we
discuss as a limitation in Section~\ref{sec:analysis}.

\section{Related Work}

\textbf{Steering Idioms (IdioSteer).} \citet{idiosteer} construct a steering vector as the
mean difference between residual-stream activations collected in figurative-context and
literal-context idiom sentences, at a fixed source layer, and inject it additively at an
earlier layer during generation (cross-layer steering, found to outperform same-layer
injection). They evaluate on a 20-idiom, ambiguous-prefix benchmark using an LLM-as-judge
(Gemma-4-31B-it, selected by benchmarking several candidate judges against a 280-item
human-annotated gold set with $\kappa=0.867$ inter-annotator agreement; the chosen judge
reaches 90.0\% accuracy and $\kappa=0.821$ against gold) and
report a sharp, coefficient-dependent shift in literal rate. Our project reuses this exact
protocol -- vector construction, injection point, coefficient grid, and judge-based
evaluation -- and asks whether it generalizes beyond the idiom class it was built on.

\textbf{IdioLink.} \citet{idiolink} introduce a retrieval benchmark linking idiomatic
spans to conceptually equivalent literal paraphrases across 107 idioms, 10{,}700
documents, and 2{,}140 queries. We use IdioLink's idiom/paraphrase pairs as the source
pool for constructing our steering vector's contrastive (figurative vs. literal)
activations, rather than for retrieval, and note it as a natural conceptual bridge between
the idiom-linking and idiom-steering lines of work.

\textbf{Activation steering more broadly.} Residual-stream activation addition
\citep{actadd} and its contrastive-prompt formalization, Contrastive Activation Addition
\citep{caa}, are the constructions IdioSteer's own mean-difference vector follows, and
which we in turn reuse for our steering vector. This line of work has otherwise mostly
targeted surface-level or coarse behavioral properties (sentiment, refusal, truthfulness,
writing style), and recent work has raised the concern that not every
semantic direction is equally steerable -- some target behaviors are not linearly
reachable from a fixed base direction, a phenomenon discussed under \emph{steerability} and
\emph{non-surjectivity} \citep{steerability}. This is directly relevant to our null result
on novel metaphors (Section~\ref{sec:analysis}): a flat effect in one direction is
consistent with a ceiling/floor effect in the data, but is also consistent with the
idiom-derived direction being only partially aligned with whatever direction would move
novel-metaphor readings.

\section{Methodology}

\subsection{Steering framework}

We follow IdioSteer's protocol directly. Given paired activations $h_i^{\text{fig}}$ and
$h_i^{\text{lit}}$ (residual stream at a source layer $L_s$, at the final token of the
idiom/expression span, for the same item under a figurative- vs. literal-biasing context),
we compute the mean-difference vector (following the paper's convention, oriented toward
the figurative class)
\[
v = \frac{1}{N}\sum_{i=1}^{N} \left(h_i^{\text{fig}} - h_i^{\text{lit}}\right), \qquad
v_{\text{unit}} = v / \lVert v \rVert.
\]
At generation time we intercept the residual stream at an earlier injection layer $L_i$
and add $\alpha \cdot v_{\text{unit}}$ at the final prompt token only (a single-token,
prompt-pass intervention -- we do not re-inject at every generated token). Positive
$\alpha$ pushes toward the figurative class, negative $\alpha$ toward literal. Following
IdioSteer's calibrated configuration for our model, we use $L_s = 14$, $L_i = 2$
(cross-layer injection), and a coefficient grid $\alpha \in
\{-4.78,-3.59,-2.39,-1.20,0,+1.20,+2.39,+3.59,+4.78\}$, with $n=5$ generated samples per
(item, $\alpha$) cell. The vector is built once, on IdioLink data
(Section~\ref{sec:idiom-data}), and reused unmodified for every experiment in this paper,
including the figurative-benchmark evaluation -- it is never retrained on figurative
data.

Model: \texttt{meta-llama/Llama-3.2-3B}, base (non-instruct) checkpoint, fp16. We use a
base rather than instruction-tuned model to match IdioSteer's setup and to keep
continuations free-form rather than shaped by chat templates. Generation uses temperature
0.7, top-$p$ 0.9, and up to 50 new tokens per completion, matching IdioSteer's reported
generation settings exactly.

\subsection{Idiom dataset}
\label{sec:idiom-data}

The steering vector's contrastive pairs are drawn from IdioLink \citep{idiolink} via its
public \texttt{Intellexus/IdioLink} release, combining the \texttt{indexes} and
\texttt{queries} configs and applying the paper's exact-surface-form filter so that each
pair genuinely contrasts a figurative and a literal use of the same idiom. This mirrors
IdioSteer's own training pool construction, which retains 1{,}683 IdioLink sentences
across 22 idioms after the same filter (a moderately imbalanced $\sim$37\% figurative /
$\sim$63\% literal split); their evaluation idioms are entirely disjoint from this
training set, and we follow the same idiom-disjoint separation between vector
construction and evaluation. For evaluation, we built a 20-idiom, ambiguous-prefix
benchmark in IdioSteer's style (IdioSteer's own benchmark is 100 sentences, 5 variants
each, over 20 idioms sourced from MAGPIE \citep{magpie}; our reconstruction is a smaller,
41-prefix-variant subset in the same style, since IdioSteer's exact released stimuli were
not available to us): each prefix ends at the idiom itself and is written so that a
reader cannot tell, from the prefix alone, whether a figurative or literal continuation is
coming (e.g. \emph{``After months of delays, the city council finally gave the project the
green light''}). Idioms include e.g. \emph{green light}, \emph{in the pink}, \emph{lie
low}, \emph{on ice}, \emph{out of the blue}, \emph{take a bow}, \emph{under the sun}.

\subsection{Non-idiomatic figurative benchmark}
\label{sec:fig-data}

We built a 100-item benchmark of non-idiomatic figurative expressions, split into three
categories: 10 conventional metaphors, 5 novel (invented) metaphors, and 5 similes, each
with 5 wording-diverse ambiguous-prefix variants (50 / 25 / 25 items respectively). Each
item is constructed to mirror the idiom benchmark's structure exactly: an ambiguous prefix
ending at the figurative expression, with both a plausible literal continuation and a
plausible figurative continuation by design. For example:

\begin{quote}
\emph{``Desperate to hide the massive embezzlement scheme before the auditors arrived, the
CFO attempted to bleach the ledger''} -- figurative continuation: \emph{aggressively wiping
the digital transaction histories from the cloud servers}; literal continuation:
\emph{pouring physical chlorine over the paper pages to dissolve the ink}.
\end{quote}

Conventional metaphors reuse widely lexicalized figurative senses (\emph{feed the fire},
\emph{build a wall}, \emph{cut the ties}, \emph{plant the seeds}); novel metaphors invent
non-lexicalized figurative extensions (\emph{bleach the ledger}, \emph{drain the
aquarium}, \emph{overwrite the DNA}) to probe whether compositionality, not just
figurativeness, matters; similes use explicit \emph{like}/\emph{as} comparisons
(\emph{like a severed cable}, \emph{like an empty vault}, \emph{like an overgrown
garden}). Each item carries \texttt{expected\_literal} and \texttt{expected\_figurative}
gold continuations used to calibrate the judge prompt.

\subsection{Evaluation protocol}

We adapt IdioSteer's LLM-as-judge protocol: given a prefix and a generated continuation,
the judge is asked to (i) state the intended figurative meaning, (ii) state the literal
meaning, (iii) check coherence, and (iv) then decide a label in
\{\textsc{figurative}, \textsc{literal}, \textsc{incoherent}\}, mirroring IdioSteer's
own judge structure and its two documented judge failure modes (not conflating a literal
\emph{physical} action with a literal reading when it happens to realize the idiom's own
conventional frame, e.g. bowing on stage for \emph{take a bow}; not over-labeling
unusual-but-coherent literal continuations as incoherent).

Our first labeling pass used a local \textsc{Qwen2.5-3B-Instruct} judge and produced a
flat literal rate across the entire idiom coefficient grid ($\approx$0.30--0.41, no clear
monotonic trend; archived in \texttt{results/labeled\_idiom\_results\_qwen3b\_baseline.csv}),
in sharp contrast to IdioSteer's reported effect at the matching configuration. Since the
steering mechanics (layers, coefficients, model) exactly match the paper, we attributed
this to insufficient judge capacity rather than a failed steering vector, and re-labeled
both the idiom and figurative results with \textsc{Qwen/Qwen3-14B} as judge. All
quantitative results in Section~\ref{sec:results} use this Qwen3-14B judge. We report,
per (dataset, category, $\alpha$) cell: figurative rate, literal rate, incoherent rate,
and coherence rate (figurative + literal).

\section{Experiments and Results}
\label{sec:results}

\subsection{Idiom baseline replication}

Table~\ref{tab:pooled} (idiom columns) shows literal rate by coefficient, pooled over all
20 idioms ($n=200$ generations per $\alpha$). Literal rate rises from 0.505 at $\alpha=0$
to 0.620 at $\alpha=-4.78$ (a $+0.115$ shift), the expected direction. Coherence stays
$\geq$0.915 in every cell, so the vector is not degrading fluency to achieve this shift.
However, this replication is markedly weaker than IdioSteer's own reported result at this
exact configuration (their baseline literal rate is $\approx$0.12, rising to
$\approx$0.49 at $\alpha=-4.78$ -- a $\sim$4$\times$ increase from a low floor). Our
baseline instead sits near 50/50, leaving less headroom for a literal push, and the
positive-$\alpha$ side is non-monotonic (e.g. $\alpha=+4.78$ also reads more literal than
baseline, 0.600 vs. 0.505), which the original paper's result does not show. We discuss
possible causes in Section~\ref{sec:analysis}.

\subsection{Applying the idiom vector to the figurative benchmark}

Table~\ref{tab:pooled} (figurative columns) shows the same idiom-derived vector applied,
unmodified, to the 100-item figurative benchmark ($n=500$ per $\alpha$). Literal rate
follows a clean, near-monotonic downward trend as $\alpha$ increases from negative to
positive: 0.774 at $\alpha=-3.59$ down to 0.652 at $\alpha=+2.39$, i.e. the vector
\emph{does} shift figurative/literal rate on data it was never built from, and in the
direction consistent with IdioSteer's own convention (negative $\alpha \to$ more literal).
Coherence again stays high throughout ($\geq$0.906).

\begin{table}[h]
\centering
\small
\begin{tabular}{@{}rrrrr@{}}
\toprule
$\alpha$ & idiom literal rate & figurative literal rate & idiom coherence & figurative coherence \\
\midrule
$-4.78$ & 0.620 & 0.754 & 0.915 & 0.918 \\
$-3.59$ & 0.605 & 0.774 & 0.940 & 0.966 \\
$-2.39$ & 0.635 & 0.770 & 0.945 & 0.938 \\
$-1.20$ & 0.580 & 0.726 & 0.970 & 0.940 \\
$0.00$  & 0.505 & 0.716 & 0.945 & 0.968 \\
$+1.20$ & 0.480 & 0.682 & 0.970 & 0.938 \\
$+2.39$ & 0.560 & 0.652 & 0.965 & 0.936 \\
$+3.59$ & 0.515 & 0.654 & 0.940 & 0.906 \\
$+4.78$ & 0.600 & 0.662 & 0.955 & 0.908 \\
\bottomrule
\end{tabular}
\caption{Literal and coherence rate by steering coefficient, pooled over idioms
($n=200$/cell) and the full figurative benchmark ($n=500$/cell). Negative $\alpha$ is the
literal-pushing direction.}
\label{tab:pooled}
\end{table}

\subsection{Breakdown by figurative category}

Table~\ref{tab:category} splits the figurative-benchmark result by category. Similes
(+0.080) and conventional metaphors (+0.040) shift in the same direction as idioms;
novel metaphors show a flat, slightly negative delta ($-0.008$) on the literal-pushing
side. Novel metaphors also have the most literal-skewed baseline of the three categories
(0.704, vs.\ 0.720 conventional / 0.720 simile\footnote{Baselines are close across
categories; the more diagnostic difference is in how much each category still moves once
steered.}) and the highest incoherent rate at extreme coefficients (0.152 at
$\alpha=+3.59$, vs.\ 0.040 at $\alpha=0$), suggesting this category is closer to a
ceiling where a further literal push has little room to register, rather than being
unaffected by steering \emph{per se}. Consistent with that reading, novel metaphors do
move in the expected direction on the figurative-pushing side: literal rate falls from
0.704 at baseline to 0.624--0.656 across $\alpha \in \{+1.20,\dots,+4.78\}$.

\begin{table}[h]
\centering
\small
\begin{tabular}{@{}lrrr@{}}
\toprule
Category & baseline literal rate ($\alpha=0$) & literal rate at $\alpha=-4.78$ & $\Delta$ \\
\midrule
Simile & 0.720 & 0.800 & $+0.080$ \\
Conventional Metaphor & 0.720 & 0.760 & $+0.040$ \\
Novel Metaphor & 0.704 & 0.696 & $-0.008$ \\
\bottomrule
\end{tabular}
\caption{Per-category literal rate shift under the maximal literal-pushing coefficient.
Full per-$\alpha$ breakdown in \texttt{analysis/summary\_by\_category\_alpha.csv}.}
\label{tab:category}
\end{table}

\begin{figure}[h]
\centering
\includegraphics[width=0.75\textwidth]{figures/literal_rate_vs_alpha.png}
\caption{Literal rate vs.\ steering coefficient $\alpha$: idioms vs.\ each figurative
category, all under the single idiom-derived steering vector.}
\label{fig:main}
\end{figure}

\section{Analysis and Discussion}
\label{sec:analysis}

\textbf{Qualitative flips.} At the item level (majority vote over 5 samples), 8 of 40
idiom (idiom, variant) pairs flip from majority-figurative at $\alpha=0$ to
majority-literal at $\alpha=-4.78$ (\emph{green light}, \emph{in the pink}, \emph{lie
low} $\times$2 variants, \emph{on ice}, \emph{out of the blue}, \emph{take a bow},
\emph{under the sun}), and 15 of 125 figurative (category, expression, variant) triples
show the same pattern, spread across all three categories. A representative figurative
flip, \emph{bleach the ledger} (novel metaphor):

\begin{quote}
\emph{``Desperate to hide the massive embezzlement scheme before the auditors arrived, the
CFO attempted to bleach the ledger...''}
\begin{itemize}[nosep]
\item Baseline ($\alpha=0$, judged \textsc{figurative}): \emph{``...entries in the general
ledger. The system had a built-in capability to reverse a journal entry, so the CFO was
able to accomplish his task...''}
\item Steered ($\alpha=-4.78$, judged \textsc{literal}): \emph{``...card stock of the check
register. However, the CFO was only able to bleach the check register for about 20
minutes before the bleach ran out...''}
\end{itemize}
\end{quote}

The steered continuation reinterprets ``bleach'' as literal chemical bleach acting on
paper, exactly the intended literal reading -- a clean example of the idiom-trained
direction successfully literalizing a metaphor it never saw during vector construction.

\textbf{Why is the idiom replication weaker than the paper's?} Three non-exclusive
explanations are consistent with our data. First, our idiom benchmark's ambiguous
prefixes may simply be less literal-biased by construction than the paper's -- our
baseline sits near 50/50 rather than their reported $\approx$0.12, so the two are not
directly comparable in absolute terms, only in direction and (weakly) in relative size.
Second, judge differences remain a live confound: IdioSteer's final judge is a much
larger model (Gemma-4-31B-it) calibrated against 280 human-annotated items with
$\kappa=0.821$ agreement to gold; our Qwen3-14B judge, while a substantial upgrade over
our first-pass Qwen2.5-3B judge (which showed no usable signal at all), has not been
calibrated against a human gold set in this project. Third, with only 20 idioms and 5
samples per cell, single-item judge noise can move pooled rates by several points, which
is consistent with the non-monotonicity we see on the positive-$\alpha$ side.

\textbf{Does the idiom vector reflect a general axis?} The evidence is mixed but leans
toward partial support. The direction of effect matches idioms for 2 of 3 non-idiomatic
categories, and the figurative benchmark as a whole shows the \emph{cleanest} monotonic
$\alpha$-vs-literal-rate curve of any dataset in this study (Figure~\ref{fig:main}),
which would be a strange coincidence if the idiom vector carried no signal at all for
non-idiomatic figurative language. At the same time, the pooled effect size on the
figurative benchmark ($+0.038$ at $\alpha=-4.78$ vs.\ baseline) is roughly a third of the
idiom effect size ($+0.115$), and one full category (novel metaphor) shows no effect in
the literal-pushing direction. This is compatible with a vector that captures a real but
\emph{partial} component of a general figurative-literal axis -- plausibly overlapping
more with the conventionalized, memorized character shared by idioms and conventional
metaphors than with the fully compositional, on-the-fly figuration required for novel
metaphors.

\textbf{Limitations.} (i) Our idiom benchmark is a hand-reconstructed 20-idiom subset,
not IdioSteer's original stimuli, so absolute rates are not directly comparable across
papers. (ii) The judge is not human-validated in this project; a small human-checked
subset (as IdioSteer uses) would strengthen confidence in the reported rates. (iii) We
did not train a second, figurative-derived steering vector to compare geometry against
the idiom vector (IdioSteer's optional extension), so we cannot say whether a
figurative-specific direction would generalize better, or how aligned the two directions
are. (iv) Sample sizes per category (5 expressions $\times$ 5 variants for novel metaphor
and simile) are small enough that category-level deltas should be read as suggestive
trends, not precise effect-size estimates.

\section{Conclusion and Future Work}
\label{sec:conclusion}

We asked whether the steering direction discovered by IdioSteer for idioms reflects a
general figurative-vs-literal axis, or is specific to idioms. Applying the idiom-derived
vector, unmodified, to a new 100-item benchmark of non-idiomatic figurative expressions,
we find partial generalization: similes and conventional metaphors shift in the expected
direction (smaller in magnitude than idioms), while novel metaphors show no effect in
that direction, plausibly due to an already literal-skewed baseline rather than true
insensitivity to steering. Our own idiom replication, while directionally consistent with
IdioSteer, is substantially weaker than their reported effect, which tempers how strongly
we can interpret the cross-category comparison. The most informative next steps are (1)
a human-validated judge subset to rule out judge noise as the source of the weak idiom
replication, and (2) training a second steering vector directly from the figurative
benchmark and measuring its cosine similarity to the idiom vector (IdioSteer's optional
Step~4, not attempted here) -- if the two vectors are highly aligned despite the idiom
vector's smaller transfer effect, that would be stronger evidence for a shared underlying
axis than the behavioral comparison alone can provide.

\bibliographystyle{plainnat}
\begin{thebibliography}{9}

\bibitem[Anonymous, 2026]{idiosteer}
Anonymous.
\emph{Steering Idioms: Controlling Figurative-vs-Literal Interpretation via
Residual-Stream Activation Steering}.
Anonymous ACL submission (under double-blind review), 2026. Provided as course
material.

\bibitem[Hashiloni et al., 2026]{idiolink}
Hashiloni et al.
\emph{IdioLink: Retrieving Meaning Beyond Words Across Idiomatic and Literal
Expressions}.
arXiv:2605.22247, 2026.

\bibitem[Turner et al., 2024]{actadd}
Turner et al.
\emph{Steering Language Models With Activation Addition}, 2024.

\bibitem[Rimsky et al., 2024]{caa}
Rimsky et al.
\emph{Steering Llama 2 via Contrastive Activation Addition}, 2024.

\bibitem[Haagsma et al., 2020]{magpie}
Haagsma et al.
\emph{MAGPIE: A Large Corpus of Potentially Idiomatic Expressions}, 2020.

\bibitem[Anonymous, 2026]{steerability}
Anonymous.
\emph{On the Steerability and Non-Surjectivity of Activation Steering
Directions}. ICLR 2026 submission. \url{https://iclr.cc/virtual/2026/10016066}

\end{thebibliography}

\end{document}
